\documentclass[10pt,a4paper]{amsart}
\usepackage[margin=19mm]{geometry}
\usepackage{amsmath,amssymb,mathtools}
\usepackage[scr=rsfs,cal=euler]{mathalfa}
\usepackage{xfrac}
\usepackage[unicode,hidelinks,bookmarks=false]{hyperref}
\newcommand{\ie}{i.e.\ }
\newcommand{\cf}{cf.\ }
\newcommand{\resp}{respectively\ }
\newcommand{\Subsection}{\subsection}
\providecommand{\nobreakdash}{\nobreak\hbox{-}}
\setlength{\emergencystretch}{2em}
\multlinegap0pt
\usepackage{bm}
\usepackage{bbm}
\usepackage{dsfont}
\usepackage{xspace}
\usepackage{cancel}
\usepackage{mathtools}
\usepackage{amsthm}
\makeatletter
\@ifundefined{theorem}{\newtheorem{theorem}{Theorem}}{}
\makeatother
\title[Dimension Of Inhomogeneous Sub-Self-Similar Sets]{Dimension Of Inhomogeneous Sub-Self-Similar Sets}
\author{Shivam Dubey, Saurabh Verma}
\date{Source: arXiv:2508.18706v3 (CC BY 4.0)}
\begin{document}
\maketitle

\noindent\emph{Source and licence.} Dimension Of Inhomogeneous Sub-Self-Similar Sets. Version arXiv:2508.18706v3 (CC BY 4.0), distributed under the Creative Commons Attribution 4.0 International licence, as recorded in the arXiv licence metadata for this exact version and shown on the abstract page https://arxiv.org/abs/2508.18706v3. Full rights notice: licenses/arxiv-2508.18706-CC-BY-4.0.txt.\par\smallskip

\noindent\textit{Editorial scope.} Counted unit: the Theorem whose source label is \texttt{None} in the article (source0.tex lines 290--296, source label \texttt{None}), delivered together with the article's own complete original proof of it (source0.tex lines 297--369, one \texttt{proof} environment of 73 lines). No mathematics is written, repaired or re-derived: statement and proof are the article's own text line for line. Required context retained uncounted: eq4 (lines 155--157). Every definition, display and label of those lines is retained unchanged. Omitted from this package and not redistributed: 1--154, 158--289, 370--928. Nothing inside a retained excerpt is omitted or altered: every retained body is delivered exactly as the article's lines are written, so no omission receipt of any kind accompanies this package. Citations: the retained excerpts carry no citation command at all, so no bibliography entry of the article is transcribed and no reference is redistributed. Reference adaptation: none was needed and none was made; every reference inside the delivered unit resolves inside it, so no unresolved reference or citation is produced. Printed slips kept literal: no printed word, symbol, hypothesis or number of the counted statement or of its proof was altered. Layout and class: the article's own class is \texttt{amsart} with packages including amssymb, srcltx, a4wide, placeins, amsmath, amsthm, amstext, cite, epsfig, float, enumerate, color, lineno, xcolor; the wrapper is the lane's pinned \texttt{amsart} editorial wrapper at 10pt on A4, which restates the theorem-like environments the retained text needs and adds the article's own macro definitions that the retained text needs (none), with extra wrapper packages (bm, bbm, dsfont, xspace, cancel, mathtools). The delivered numbering is the unit's own. Assets: no retained span contains a figure environment, a graphics-inclusion command, a PDF-page inclusion, a table or a TikZ picture; no image file is copied into this package, the package declares no asset dependency and no image file is redistributed with it. Licence: version arXiv:2508.18706v3 of this article is distributed under the Creative Commons Attribution 4.0 International licence, as recorded in the arXiv licence metadata for this exact version and shown on the abstract page https://arxiv.org/abs/2508.18706v3. A full rights notice is delivered at licenses/arxiv-2508.18706-CC-BY-4.0.txt. Characters: every retained excerpt is pure ASCII, so the delivered engine, XeLaTeX with its default Latin Modern text font, reports no missing character.
\begin{equation}\label{eq4}
F_C \subseteq \bigcup_{i=1}^N f_i(F_C) \cup C.
\end{equation}

\begin{theorem}
Let $E$ be an SSS set associated with the IFS $\mathcal{I}$ such that $\Theta(S) = E$, and let $C \subset X$ be a non-empty compact set. Define
\[
O_S := \bigcup_{{\omega} \in S^*} f_{{\omega}}(C).
\]
Then $ E\cup O_S$ is an ISSS set corresponding to the IFS $\mathcal{I}$ with condensation set $C$.
\end{theorem}

\begin{proof}
To prove that $E\cup O_S$ is an ISSS set, we first show that it is compact. Since $O_S$ is bounded, it suffices to prove that $E \cup O_S$ is closed. Let $\{x_n\}$ be a sequence in $E \cup O_S$ converging to some $x \in X$. If $x_n \in E$ for infinitely many $n$, then there exists a subsequence $\{x_{n_r}\}$ with $x_{n_r} \in E$. Since $E$ is closed, $\lim x_{n_r} = x \in E \subseteq E \cup O_S$. If $x_n \in E$ for only finitely many $n$, i.e., $\{x_n\}$ is eventually in $O_S$, then there exists $n_0$ such that $x_n \in O_S$ for all $n \geq n_0$. Hence, for each $n \geq n_0$, there exists a string $\omega_{|_n} \in S^*$ such that $x_n \in f_{\omega_{|_n}}(C)$.
Let us denote the length of any string $\omega := \omega_1 \omega_2 \dots \omega_k \in S^*$ by $|\omega|=k.$

\textbf{Case 1:} $\liminf |\omega_{|_n}| = k < \infty$. Then there exists a subsequence $\{\omega_{{|_n}_j}\}$ of $\{\omega_{|_n}\}$ such that $\omega_{{|_n}_j} \in S^k$ for all $j$, where $S^k$ is a finite set. Thus, there exists a string $t \in \{\omega_{{|_n}_j} \mid j \in \mathbb{N}\}$ and a further subsequence $\{\omega_{{|_n}_{j_p}}\}$ such that $\omega_{{|_n}_{j_p}} = t$ for all $p$. Then $x_{n_{j_p}} \in f_t(C)$, and since $f_t(C)$ is closed, $\lim x_{n_{j_p}} = x \in f_t(C) \subseteq O_S$.

\textbf{Case 2:} $\liminf |\omega_{|_n}| = \infty$. Then
\[
\begin{aligned}
\operatorname{diam}(f_{\omega_{|_n}}(C) \cup f_{\omega_{|_n}}(E)) &= \rho_{\omega_{|_n}} \operatorname{diam}(C \cup E) \\
&\leq \rho_{\max}^{|\omega_{|_n}|} \operatorname{diam}(C \cup E) \to 0 \quad \text{as } n \to \infty,
\end{aligned}
\]
since $\rho_{\max} = \max\{\rho_i\} < 1$.

Now, for each $n$, choose $y_n \in f_{\omega_{|_n}}(E)$. For positive integers $m$ and $n$,
\[
\begin{aligned}
|y_n - y_m| &\leq |y_n - x_n| + |x_n - x_m| + |x_m - y_m| \\
&\leq \operatorname{diam}(f_{\omega_{|_n}}(C) \cup f_{\omega_{|_n}}(E)) + |x_n - x_m| + \operatorname{diam}(f_{\omega_{|_m}}(C) \cup f_{\omega_{|_m}}(E)).
\end{aligned}
\]
Hence, $\{y_n\}$ is a Cauchy sequence and converges to some $y \in X$. Then,
\[
\begin{aligned}
|x - y| &\leq |x - x_n| + |x_n - y_n| + |y_n - y| \\
&\leq |x - x_n| + \operatorname{diam}(f_{\omega_{|_n}}(C) \cup f_{\omega_{|_n}}(E)) + |y_n - y| \to 0.
\end{aligned}
\]
Thus, $x = y$.

Since $y_n \in f_{\omega_{|_n}}(E)$ for each $n$, there exists $y' \in E$ such that $y_n = f_{\omega_1 \ldots \omega_n}(y')$. Let $\epsilon > 0$. Choose $n_1 \in \mathbb{N}$ such that for all $n \geq n_1$,
\[
|y_n - y| = |f_{\omega_1 \ldots \omega_n}(y') - y| < \epsilon / 2.
\]
Since $\Theta(S) = E$, we have $\lim_{n \to \infty} f_{\omega_1 \ldots \omega_n}(y') = z_{n_\epsilon} \in E$. Choose $n_2 \in \mathbb{N}$ such that for all $n \geq n_2$,
\[
|f_{\omega_1 \ldots \omega_n}(y') - z_{n_\epsilon}| < \epsilon / 2.
\]
Let $n' = \max\{n_1, n_2\}$. Then,
\[
\begin{aligned}
|y - z_{n_\epsilon}| &\leq |f_{\omega_1 \ldots \omega_n}(y') - y| + |f_{\omega_1 \ldots \omega_n}(y') - z_{n_\epsilon}| \\
&< \epsilon / 2 + \epsilon / 2 = \epsilon.
\end{aligned}
\]
Since $E$ is closed, $x = y = \lim z_{n_\epsilon} \in E \subseteq E \cup O_S$. Thus, $E \cup O_S$ is closed, and hence compact.

It remains to show that $E \cup O_S$ satisfies Equation~\eqref{eq4}.

Clearly,
\[
\begin{aligned}
E \cup O_S &\subseteq \bigcup_{i=1}^N f_i(E) \cup O_S \\
&\subseteq \bigcup_{i=1}^N f_i(E \cup O_S) \cup O_S \cup C.
\end{aligned}
\]

Now we show that $O_S \subseteq \bigcup_{i=1}^N f_i(O_S)$. Let $x \in O_S$. Then there exists $\omega_{|_n} \in S^*$ and $c \in C$ such that
\[
x = f_{\omega_{|_n}}(c) = f_{\omega_1 \ldots \omega_n}(c).
\]
Let $x' = f_{\omega_2 \ldots \omega_n}(c) \in O_S$. Then,
\[
x = f_{\omega_1}(x') \in \bigcup_{i=1}^N f_i(O_S).
\]

Thus,
\[
E \cup O_S \subseteq \bigcup_{i=1}^N f_i(E \cup O_S) \cup C.
\]
This completes the proof.
\end{proof}

\end{document}
